\section{Space Complexity --- Codec --- Valid}
%%% generated by the blueprint pipeline from TCSlib.Complexity.SpaceComplexity.Codec.Valid
\subsection*{Space Complexity --- Codec --- Valid: Introduction}

Declarations of the Lean module \texttt{TCSlib.Complexity.SpaceComplexity.Codec.Valid}.

\subsection*{Space Complexity --- Codec --- Valid: Results}

\begin{definition}[Partial chunk decoder]
\lean{Complexity.Codec.decChunk}
\label{Complexity.Codec.decChunk}
\leanok
\uses{Complexity.Codec.CellV, Complexity.Codec.encChunk}
Let $N$ be a nondeterministic Turing machine with Boolean input alphabet, and let $ch$
be a finite Boolean list. The decoder returns a cell value whose encoding is $ch$ when
such a value exists, and returns no value when none exists.
\end{definition}

\begin{lemma}[Chunk decoding recovers its encoding]
\lean{Complexity.Codec.decChunk_encChunk}
\label{Complexity.Codec.decChunk_encChunk}
\leanok
\uses{Complexity.Codec.CellV, Complexity.Codec.decChunk, Complexity.Codec.encChunk, Complexity.Codec.encChunk_injective}
\difficulty{4}
For every machine $N$ and every cell value $v$, decoding the encoded chunk of $v$
returns $v$.
\end{lemma}

\begin{lemma}[Encoding after decoding]
\lean{Complexity.Codec.encChunk_of_decChunk}
\label{Complexity.Codec.encChunk_of_decChunk}
\leanok
\uses{Complexity.Codec.CellV, Complexity.Codec.decChunk, Complexity.Codec.encChunk}
\difficulty{4}
Let $N$ be a finite nondeterministic Turing machine with Boolean tape symbols, and let
$\mathrm{CellV}$ be its cell-value type. If decoding a Boolean list $ch$ yields the cell
value $cv$, then encoding $cv$ reproduces $ch$.
\end{lemma}

\begin{definition}[Verifier state space]
\lean{Complexity.Codec.VSt}
\label{Complexity.Codec.VSt}
\leanok
\uses{Complexity.Codec.Hd}
For a nondeterministic Turing machine $N$ with Boolean input alphabet, a verifier state
consists of a Boolean flag, an optional pair comprising an optional machine state and
one of three head directions, a Boolean-valued function on the machine's $k$ tape
indices, an optional-Boolean-valued function on those indices, and two further Boolean
flags.
\end{definition}

\begin{definition}[Coordinated blank cell values]
\lean{Complexity.Codec.BlankTogether}
\label{Complexity.Codec.BlankTogether}
\leanok
\uses{Complexity.Codec.CellV}
Let $N$ be a nondeterministic Turing machine with Boolean input alphabet, and let $cv$
be a cell value. Its tape-entry components are all absent, and its Boolean components
are either all true or all false.
\end{definition}

\begin{definition}[Initial verifier state]
\lean{Complexity.Codec.vInit}
\label{Complexity.Codec.vInit}
\leanok
\uses{Complexity.Codec.VSt}
For a nondeterministic Turing machine $N$ with Boolean input alphabet, the initial
verifier state has first flag true, no stored head state, all tape-index Boolean values
false, all optional-Boolean values absent, and both final flags true.
\end{definition}

\begin{definition}[Verifier state update]
\lean{Complexity.Codec.vStep}
\label{Complexity.Codec.vStep}
\leanok
\uses{Complexity.Codec.BlankTogether, Complexity.Codec.CellV, Complexity.Codec.VSt}
Let $N$ be a nondeterministic Turing machine with Boolean input alphabet, $st$ a
verifier state, and $cv$ a cell value. The updated state has first flag false; its
stored head state is $cv$'s when the old first flag is true and otherwise is retained.
At each tape index, the Boolean value becomes the disjunction of its old value and the
corresponding Boolean component of $cv$, while the optional-Boolean value becomes the
corresponding component of $cv$ when that Boolean component is true and is otherwise
retained. The next-to-last flag remains true only if it was previously true, the head
state is present exactly when required by the old first flag, and no index has both its
old Boolean value and its component of $cv$ true; the last flag remains true only if it
was previously true and $cv$ has coordinated blank components.
\end{definition}

\begin{definition}[Optional verifier state update]
\lean{Complexity.Codec.vStepO}
\label{Complexity.Codec.vStepO}
\leanok
\uses{Complexity.Codec.CellV, Complexity.Codec.VSt, Complexity.Codec.vStep}
Let $N$ be a nondeterministic Turing machine with Boolean input alphabet and $st$ a
verifier state. On no cell value, the update retains the first flag, stored head state,
and both families of tape-index values; it sets the next-to-last flag to false and
retains the last flag. On a supplied cell value, it applies the verifier state update
defined by that value.
\end{definition}

\begin{definition}[Verifier update from a chunk]
\lean{Complexity.Codec.vChunk}
\label{Complexity.Codec.vChunk}
\leanok
\uses{Complexity.Codec.VSt, Complexity.Codec.decChunk, Complexity.Codec.vStepO}
Let $N$ be a nondeterministic Turing machine with Boolean input alphabet, $st$ a
verifier state, and $ch$ a finite list of pairs of Booleans. The update decodes the list
of first components of $ch$ as a cell value and applies the optional verifier state
update to $st$ using that decoded value.
\end{definition}

\begin{definition}[Verifier scan state]
\lean{Complexity.Codec.vScan}
\label{Complexity.Codec.vScan}
\leanok
\uses{Complexity.Codec.VSt, Complexity.Codec.adStep, Complexity.Codec.hH, Complexity.Codec.vChunk, Complexity.Codec.vInit}
Let $N$ be a finite nondeterministic Turing machine with Boolean input symbols, and let
$u$ be a finite Boolean string. The scan state is the state component obtained by
starting from the prescribed initial scan data and folding the verifier update operation
over the list of pairs $(b,b)$ for $b$ in $u$, in order.
\end{definition}

\begin{definition}[Verifier acceptance condition]
\lean{Complexity.Codec.VAcc}
\label{Complexity.Codec.VAcc}
\leanok
\uses{Complexity.Codec.VSt}
Let $N$ be a finite nondeterministic Turing machine with Boolean input symbols, and let
$st$ be a verifier state. The state is accepting exactly when its designated Boolean
flag is true and every Boolean entry of its work-tape flag vector is true.
\end{definition}

\begin{theorem}[Summary of a finite sequence of cell updates]
\lean{Complexity.Codec.foldl_vStep}
\label{Complexity.Codec.foldl_vStep}
\leanok
\uses{Complexity.Codec.BlankTogether, Complexity.Codec.CellV, Complexity.Codec.vInit, Complexity.Codec.vStep}
\difficulty{6}
Let $N$ be a finite nondeterministic Turing machine with Boolean tape symbols, and let
$f:\mathbb{N}\to\mathrm{CellV}(N)$ and $J\in\mathbb{N}$. Starting from the initial
validation state and successively applying the cell update to $f(0),\ldots,f(J-1)$, the
first state flag is true exactly when $J=0$, and the stored head data is absent when
$J=0$ and otherwise equals the head data of $f(0)$. For each work tape $i$, its visited
flag is true exactly when some processed cell marks that tape; its stored symbol is
absent if no processed cell marks it, and otherwise is the symbol at the last processed
cell that marks it. The first validity flag is true exactly when among the processed
cells the head data is present precisely at index $0$, and each tape is marked at most
once; the second validity flag is true exactly when every processed cell satisfies the
blank-together condition.
\end{theorem}

\begin{lemma}[Persistence of a False Validity Flag]
\lean{Complexity.Codec.vStepO_ok_false}
\label{Complexity.Codec.vStepO_ok_false}
\leanok
\uses{Complexity.Codec.CellV, Complexity.Codec.VSt, Complexity.Codec.vStep, Complexity.Codec.vStepO}
\difficulty{3}
Let $N$ be a finite nondeterministic Turing machine with Boolean tape symbols, let $st$
be a verifier state, and let $o$ be either no cell or a cell value. If the validity flag
of $st$ is false, then the validity flag after applying the option-valued verifier step
to $st$ and $o$ is also false.
\end{lemma}

\begin{lemma}[A false validity flag persists]
\lean{Complexity.Codec.foldl_vStepO_ok_false}
\label{Complexity.Codec.foldl_vStepO_ok_false}
\leanok
\uses{Complexity.Codec.CellV, Complexity.Codec.VSt, Complexity.Codec.vStepO, Complexity.Codec.vStepO_ok_false}
\difficulty{4}
Let $N$ be a finite nondeterministic Turing machine with Boolean symbols, let $L$ be a
finite list of optional cells, and let $st$ be a validation state whose first validity
flag is false. After applying the option-valued cell update successively to every entry
of $L$, that flag remains false.
\end{lemma}

\begin{lemma}[Optional and ordinary cell updates agree]
\lean{Complexity.Codec.foldl_vStepO_some}
\label{Complexity.Codec.foldl_vStepO_some}
\leanok
\uses{Complexity.Codec.CellV, Complexity.Codec.VSt, Complexity.Codec.vStep, Complexity.Codec.vStepO}
\difficulty{3}
Let $N$ be a finite nondeterministic Turing machine with Boolean symbols, let $L$ be a
finite list of cells, and let $st$ be a validation state. Folding the option-valued
update over the list obtained by making every entry of $L$ present gives the same state
as folding the ordinary cell update over $L$.
\end{lemma}

\begin{lemma}[Flattening equal-length blocks]
\lean{Complexity.Codec.eq_flatten_blocks}
\label{Complexity.Codec.eq_flatten_blocks}
\leanok
\difficulty{4}
Let $u$ be a list over a type $\alpha$, and let $w,m$ be natural numbers. If $u$ has
length $wm$, then $u$ is the concatenation, in order, of the $m$ consecutive blocks of
length $w$ obtained by taking the first $w$ entries after dropping the first $wj$
entries, for $j=0,\ldots,m-1$.
\end{lemma}

\begin{lemma}[Finite function list and range map]
\lean{Complexity.Codec.ofFn_eq_map_range}
\label{Complexity.Codec.ofFn_eq_map_range}
\leanok
\difficulty{3}
For any type $\beta$, function $g:\mathbb{N}\to\beta$, and $m\in\mathbb{N}$, the list of
values $g(0),g(1),\ldots,g(m-1)$ is the image under $g$ of the list $0,1,\ldots,m-1$.
\end{lemma}

\begin{theorem}[Scanning a Word by Fixed-Width Blocks]
\lean{Complexity.Codec.vScan_eq}
\label{Complexity.Codec.vScan_eq}
\leanok
\uses{Complexity.Codec.W, Complexity.Codec.decChunk, Complexity.Codec.eq_flatten_blocks, Complexity.Codec.foldl_adStep_flatten, Complexity.Codec.hH, Complexity.Codec.ofFn_eq_map_range, Complexity.Codec.vChunk, Complexity.Codec.vInit, Complexity.Codec.vScan, Complexity.Codec.vStepO, Complexity.Codec.zip_flatten}
\difficulty{6}
Let $N$ be a finite nondeterministic Turing machine with Boolean tape symbols, let $u$
be a finite Boolean word, and let $m\in\mathbb N$ satisfy $|u|=Wm$, where $W$ is the
block width associated to $N$. Scanning $u$ produces the verifier state obtained by
starting at the initial state and successively applying the option-valued verifier step
to the decoded blocks $(u[Wj:W(j+1)])$ for $j=0,\ldots,m-1$.
\end{theorem}

\begin{definition}[Output summary encoding]
\lean{Complexity.Codec.outOf}
\label{Complexity.Codec.outOf}
\leanok
Let $N$ be a finite nondeterministic Turing machine with Boolean input symbols. Encode
the three output summaries by elements of $	ext{Fin}(3)$: the empty summary maps to $0$,
the accepting summary maps to $1$, and the dead summary maps to $2$.
\end{definition}

\begin{lemma}[Decoding output summary indices]
\lean{Complexity.Codec.sumIdx_outOf}
\label{Complexity.Codec.sumIdx_outOf}
\leanok
\uses{Complexity.Codec.outOf, Complexity.Codec.sumIdx, Turing.outSummary}
\difficulty{4}
For every $o\in\{0,1,2\}$, first decode $o$ as the empty list when $o=0$, the one-bit
list $[1]$ when $o=1$, and the one-bit list $[0]$ when $o=2$. The index of the resulting
output summary is $o$.
\end{lemma}

\begin{definition}[Fixed-width tape block]
\lean{Complexity.Codec.block}
\label{Complexity.Codec.block}
\leanok
\uses{Complexity.Codec.W}
Let $N$ be a finite nondeterministic Turing machine with Boolean input symbols, let $u$
be a finite Boolean string, and let $j$ be a natural number. The block at position $j$
is the length-$W(N)$ segment of $u$ beginning at index $W(N)j$, truncated if $u$ ends
before the segment is complete.
\end{definition}

\begin{lemma}[Blocks of a concatenated chunk word]
\lean{Complexity.Codec.block_chunksWord}
\label{Complexity.Codec.block_chunksWord}
\leanok
\uses{Complexity.Codec.CellV, Complexity.Codec.W, Complexity.Codec.block, Complexity.Codec.chunksWord, Complexity.Codec.encChunk, Complexity.Codec.take_drop_flatten}
\difficulty{3}
Let $N$ be a finite nondeterministic Turing machine over Boolean symbols, let $f$ assign
a cell value to each natural number, and let $m,j$ be natural numbers with $j<m$. The
block at position $j$ of the word formed by concatenating the first $m$ encoded cell
values is the encoded cell value $f(j)$.
\end{lemma}

\begin{definition}[Configuration decoded from cells]
\lean{Complexity.Codec.mkCfg}
\label{Complexity.Codec.mkCfg}
\leanok
\uses{Complexity.Codec.CellV, Complexity.Codec.outOf, Turing.Cfg}
Let $N$ be a finite nondeterministic Turing machine with Boolean input symbols, let $x$
be a finite Boolean input, let $s$ be a natural number, let $f$ assign a cell value to
each natural index, and let $pos$ be one of the $|x|+2$ input-head positions. The
decoded configuration has internal state and output summary read from cell $f(0)$, and
input-head position $pos$. On each work tape, cells with integer index $z$ satisfying
$|z|\le s$ receive the optional symbol stored at index $z+s$ of $f$; all other cells are
blank. Each work-tape head is placed at the least index $j$ with $0\le j<2s+1$ whose
corresponding flag is true, shifted by $-s$, or at $0$ if there is no such index.
\end{definition}

\begin{theorem}[Reconstruction from window cells]
\lean{Complexity.Codec.cellOf_mkCfg}
\label{Complexity.Codec.cellOf_mkCfg}
\leanok
\uses{Complexity.Codec.CellV, Complexity.Codec.cellOf, Complexity.Codec.mkCfg, Turing.Cfg.InWindow}
\difficulty{6}
Let $N$ be a finite nondeterministic Turing machine with Boolean symbols, let $x$ be a
finite Boolean string, let $s$ be a natural number, let $f$ assign a cell value to each
natural-number index, and let $p$ be a valid input-head position among the $|x|+2$
positions. Suppose among indices $j<2s+1$, the head-and-summary slot is present exactly
at $j=0$; for each work tape, at most one such index marks its head; and for each work
tape, at least one such index marks its head. Then the configuration reconstructed from
$f$, $s$, and $p$ has all work-tape heads within distance $s$ of the origin and blank
cells outside that window, and its encoded cell at every index $j<2s+1$ is exactly
$f(j)$.
\end{theorem}

\begin{definition}[Validity of a certificate string]
\lean{Complexity.Codec.ValidChk}
\label{Complexity.Codec.ValidChk}
\leanok
\uses{Complexity.Codec.VAcc, Complexity.Codec.W, Complexity.Codec.vScan, Complexity.Codec.wP, Complexity.LogProg.bitsVal}
Let $N$ be a finite nondeterministic Turing machine with Boolean input symbols, and let
$n,s$ be natural numbers and $w$ a finite Boolean string. The string is valid exactly
when its length is $W(N)(2s+1)+wP(n)+1$, the scan of its prefix of length $W(N)(2s+1)$
is accepting, the binary value of the next $wP(n)$ bits is at most $n+1$, and the bit at
index $W(N)(2s+1)+wP(n)$ is true (with out-of-range lookup interpreted as false).
\end{definition}

\begin{lemma}[Chunk prefix of the configuration word]
\lean{Complexity.Codec.cfgWord_take}
\label{Complexity.Codec.cfgWord_take}
\leanok
\uses{Complexity.Codec.W, Complexity.Codec.cellOf, Complexity.Codec.cfgWord, Complexity.Codec.chunksWord, Complexity.Codec.length_chunksWord, Turing.Cfg}
\difficulty{3}
Let $N$ be a finite nondeterministic Turing machine with Boolean tape symbols, let $x$
be a Boolean input word, let $C$ be a configuration of $N$ on input $x$, and let
$s\in\mathbb{N}$. The first $W(2s+1)$ bits of the encoded configuration word are exactly
the concatenation of the encoded chunks for the cells of that configuration.
\end{lemma}

\begin{lemma}[Input-position segment of the configuration word]
\lean{Complexity.Codec.cfgWord_pos}
\label{Complexity.Codec.cfgWord_pos}
\leanok
\uses{Complexity.Codec.W, Complexity.Codec.cfgWord, Complexity.Codec.length_bits_pos_le, Complexity.Codec.length_chunksWord, Complexity.Codec.length_padBits, Complexity.Codec.padBits, Complexity.Codec.wP, Turing.Cfg}
\difficulty{3}
Let $N$ be a finite nondeterministic Turing machine with Boolean tape symbols, let $x$
be a Boolean input word, let $C$ be a configuration of $N$ on input $x$, and let
$s\in\mathbb{N}$. The segment of the encoded configuration word obtained by dropping its
first $W(2s+1)$ bits and then taking the first $w_P(|x|)$ bits is exactly the padded bit
representation of the input-head position of $C$.
\end{lemma}

\begin{lemma}[Final bit of the configuration word]
\lean{Complexity.Codec.cfgWord_last}
\label{Complexity.Codec.cfgWord_last}
\leanok
\uses{Complexity.Codec.W, Complexity.Codec.cellOf, Complexity.Codec.cfgWord, Complexity.Codec.chunksWord, Complexity.Codec.length_bits_pos_le, Complexity.Codec.length_chunksWord, Complexity.Codec.length_padBits, Complexity.Codec.padBits, Complexity.Codec.wP, Turing.Cfg}
\difficulty{4}
Let $N$ be a finite nondeterministic Turing machine with Boolean tape symbols, let $x$
be a Boolean input word, let $C$ be a configuration of $N$ on input $x$, and let
$s\in\mathbb{N}$. The bit of the encoded configuration word at index $W(2s+1)+w_P(|x|)$
is true, where $W$ is the chunk width and $w_P(|x|)$ is the padded input-position width.
\end{lemma}

\begin{lemma}[Scanning an Encoded Configuration]
\lean{Complexity.Codec.vScan_cfgWord_eq}
\label{Complexity.Codec.vScan_cfgWord_eq}
\leanok
\uses{Complexity.Codec.W, Complexity.Codec.block, Complexity.Codec.block_chunksWord, Complexity.Codec.cellOf, Complexity.Codec.cfgWord, Complexity.Codec.cfgWord_take, Complexity.Codec.decChunk, Complexity.Codec.foldl_vStepO_some, Complexity.Codec.length_cfgWord, Complexity.Codec.vInit, Complexity.Codec.vScan, Complexity.Codec.vScan_eq, Complexity.Codec.vStep, Turing.Cfg}
\difficulty{4}
Let $N$ be a finite nondeterministic Turing machine with Boolean tape symbols, let $x$
be a finite Boolean input, let $s\in\mathbb N$, and let $C$ be a configuration of $N$ on
input $x$. Scanning the prefix of the encoded configuration word consisting of its
$2s+1$ cell blocks yields the same verifier state as starting from the initial verifier
state and processing, in order, the cells indexed by $0,1,\ldots,2s$.
\end{lemma}

\begin{theorem}[Scanning a windowed configuration encoding]
\lean{Complexity.Codec.vScan_cfgWord}
\label{Complexity.Codec.vScan_cfgWord}
\leanok
\uses{Complexity.Codec.BlankTogether, Complexity.Codec.VAcc, Complexity.Codec.W, Complexity.Codec.cellOf, Complexity.Codec.cfgWord, Complexity.Codec.foldl_vStep, Complexity.Codec.sumIdx, Complexity.Codec.vInit, Complexity.Codec.vScan, Complexity.Codec.vScan_cfgWord_eq, Complexity.Codec.vStep, Turing.Cfg, Turing.Cfg.InWindow, Turing.outSummary}
\difficulty{5}
Let $N$ be a finite nondeterministic Turing machine, let $x$ be a Boolean input, let $s$
be a natural number, and let $C$ be a configuration of $N$ on input $x$. Suppose every
work-tape head of $C$ lies within distance $s$ of the origin and every work-tape cell
outside that window is blank. Scan the initial $W(2s+1)$ bits of the encoded
configuration word by the virtual scanner. Its recorded head and output summary is
exactly the configuration's state and output summary; the scanner's acceptance condition
holds, its recorded work-tape symbols at each tape head agree with $C$, and its
blankness flag is true exactly when every window cell has the specified joint-blank
property.
\end{theorem}

\begin{lemma}[Missing cell clears the validity flag]
\lean{Complexity.Codec.foldl_vStepO_none}
\label{Complexity.Codec.foldl_vStepO_none}
\leanok
\uses{Complexity.Codec.CellV, Complexity.Codec.vInit, Complexity.Codec.vStepO, Complexity.Codec.vStepO_ok_false}
\difficulty{4}
Let $N$ be a finite nondeterministic Turing machine with Boolean symbols, let
$d:\mathbb{N}\to\mathrm{Option}({\mathrm{CellV}(N)})$, and let $m\in\mathbb{N}$. If
$d(j)$ is absent for some $j<m$, then after folding the option-valued cell update over
$d(0),\ldots,d(m-1)$ from the initial validation state, the first validity flag is
false.
\end{lemma}

\begin{theorem}[Characterization of valid configuration encodings]
\lean{Complexity.Codec.validChk_iff}
\label{Complexity.Codec.validChk_iff}
\leanok
\uses{Complexity.Codec.CellV, Complexity.Codec.ValidChk, Complexity.Codec.W, Complexity.Codec.bitsVal_padBits, Complexity.Codec.block, Complexity.Codec.cellOf, Complexity.Codec.cellOf_mkCfg, Complexity.Codec.cfgWord, Complexity.Codec.cfgWord_last, Complexity.Codec.cfgWord_pos, Complexity.Codec.chunksWord, Complexity.Codec.decChunk, Complexity.Codec.encChunk_of_decChunk, Complexity.Codec.eq_flatten_blocks, Complexity.Codec.eq_padBits, Complexity.Codec.foldl_vStep, Complexity.Codec.foldl_vStepO_none, Complexity.Codec.foldl_vStepO_some, Complexity.Codec.length_cfgWord, Complexity.Codec.mkCfg, Complexity.Codec.padBits, Complexity.Codec.vScan_cfgWord, Complexity.Codec.vScan_eq, Complexity.Codec.wP, Complexity.LogProg.bitsVal, Turing.Cfg, Turing.Cfg.InWindow}
\difficulty{6}
Let $N$ be a finite nondeterministic Turing machine with Boolean symbols, let $x$ be a
finite Boolean input, and let $s$ be a nonnegative integer. A Boolean word $w$ satisfies
the validity check for $N$, $x$, and $s$ if and only if there exists a configuration $C$
of $N$ on input $x$ such that every work-tape head lies within distance $s$ of the
origin, every work tape is blank at positions whose distance from the origin exceeds
$s$, and $w$ is the encoding of $C$.
\end{theorem}
